% !TeX root = ../papers/nc-p3.tex
\section{有限窓の一組の共通三平面と反対対称性}
\label{sec:marked-common-triple-symmetry}
\begin{theorem}\label{thm:marked-common-triple-symmetry}
単純な局所点を持つ有限Frobenius点窓で、既出の中支持3による認証がないとする。
印付き交換子表示で $q=\dim\langle Q,P\rangle=3$ または $g=\dim\langle G,F\rangle=3$ なら、
反対階数の仮定を置かずに少なくとも一方の反対位置が対称であり、$\rho\le10$ である。
$q=3$, $g>3$ なら基底を選んで $P=Q$、$G,F$ はすべて対称とでき、
$g=4,5$ に限る。交叉交換子の次元が2の場合には $g=5$ である。
\end{theorem}
\begin{proof}[交換子の二平面と非対称な基底変換]
巡回して $q=3$ とし、$P_d=Q_{Jd}$ と書く。$J$ は可逆である。
$N=\langle[G_a,F_c]\rangle$ は係数添字 $b,d$ の行列として
$KJ$ ($K$ は交代行列) の包に入る。次元 $r$ は $Q$ の三平面の交換子次元であり、
既約性とSchurの補題により $r=2,3$ である。$r=3$ なら跡零条件
$\operatorname{tr}(KJ)=0$ が全 $K$ に成り立って $J$ は対称である。
また $g=3$ の全場合は定理\ref{thm:rank-nine-finite-symmetry}の証明の共通三平面の場合による。
そこでは混合型も一方が対称、両交代型は点の有限性に反し、反対階数を用いる必要はない。

$g>3,r=2$ として $J$ が非対称なら矛盾することを示す。$D=J-J^t$ と置く。
$N$ を書く二平面の $K$ は全て $J$ と跡直交なので、その交代空間内の直交補は $kD$ である。
一般 $X\in\langle G\rangle$ で $[X,\langle F\rangle]=N$ となる。
そうでなければ $X$ に線型な $\operatorname{Hom}(\langle F\rangle,N)$ の部分空間が全て階数1以下となり、
共通像なら $r=2$ に反し、共通余像なら非零の跡零 $F$ が全 $G$ と交換してしまう。
従って $X,X^2$ は $N$ と跡直交する。多項式性から全 $X$ で同じである。$F$ 側でも同様である。

$JX$ の交代部分は $D$ の倍なので $X'=X-\lambda(X)I$ を $JX'$ が対称になるよう一意に取れる。
$J(X')^2$ の交代部分も $D$ の倍であり、$JX'=X'^tJ^t$ より
\[
                       X'^t D X'\in kD .
\]
$kr_0=\ker D$、$\pi$ をその商への射影とする。右の倍率が一般に非零なら
$\pi X'$ の二次外積と $\pi$ の二次外積が比例し、階数2の両写像の核は同じである。
従って $X'r_0\in kr_0$ が全族に成り立って既約性に反する。
よって $\pi X'$ はいつも階数1以下である。共通像の場合は全ての $X'$ が固定二平面に値を持ち、
元の $X$ もその平面を保ってしまう。残るのは共通余像 $k\alpha$ の場合で、
$X'(\ker\alpha)\subset kr_0$ である。

$\beta=Jr_0$ とすると対称形式 $JX'$ は $\ker\alpha$ を $k\beta$ に送る。
$\alpha,\beta$ が独立ならこうした対称形式は $\alpha^2,\beta^2$ の包に限られる。
実際 $\ker\alpha$ での対称性から制限は $\beta$ に $\beta$ の制限の定数倍を掛けたもので、
それを引くと両因子が $k\alpha$ に入る。$X\mapsto X'$ は跡零空間上で単射なので二次元では足りない。
従って $\alpha$ は $\beta$ に比例し、三次元の対称形式空間は強制的に
$\{\beta\otimes l+l\otimes\beta:l\in(k^3)^*\}$ 全体である。
$F$ 側でも同じ $D,J,r_0,\beta$ を使うため同じ空間となる。
$J^{-1}$ を掛けて跡零部分を取れば二つの元の三平面が等しくなり $g>3$ に反する。
これで $J$ は対称なので合同で単位行列に取れる。
\end{proof}
\begin{proof}[異なる三平面の交叉交換子]
今 $[G_a,F_c]$ は全て交代である。各三平面を跡零対称部分と外積行列に分けて
$G=(A,[u]_\times), F=(B,[v]_\times)$ と書き、
$D_u B=[[u]_\times,B]$ とすると条件は $D_u B=D_v A$ である。
どちらかの交代射影が零なら、相手にも交代部分はない。
非零外積行列と交換する跡零対称行列は一次元しかないからである。
どちらも零でなく片側の交代射影が一次元以下なら、その純対称な核の二元が相手の非零外積と
交換するためやはり不可能である。両側とも交代射影が二次元以上なら、相手の純対称核は零である。
独立な二つの外積と交換する跡零対称行列は零であるためである。
従って $A=A(u), B=B(v)$ と両者を全交代三空間上のグラフに書ける。

和の写像 $H=A+B$ は $D_u H(v)=D_v H(u)$ を満たす。
この十五元の線型系は定理\ref{thm:rank-eight-finite-symmetry}の行列式 $-3456$ の検算から核を持たない。
差側は従って $D_u A(u)=0$ を満たす。一般の $u$ では中心化空間により
$A(u)$ は $u u^t-(u^tu)I/3$ に比例する。この二次行列の成分は互いに全体公約因子を持たない
（既に三つの非対角単項式がそうである）。有理倍率で全成分が一次多項式になることは
零以外になく、$A=B=0$ である。この両交代型なら $g=3$ なので除かれる。
従って $g>3$ では両方とも跡零対称で $g\le5$ である。
星印成分と純成分を交換子表示で比較すれば $b,d$ 交換で $w$ が不変であり、$\rho\le10$ となる。

最後に $r=2$ なら交換子と跡直交する非零外積行列 $K$ を取る。
対称跡零の五空間の交代形式 $\omega(X,Y)=\operatorname{tr}(K[X,Y])$ の根基は一次元である。
実際、対称跡零空間上の跡対合は完全で、根基は $K$ と交換する跡零対称行列である。
二つの三平面は $\omega$ で互いに直交するので根基をともに含み、商の四次元シンプレクティック空間で
互いに直交補の二平面となる。二平面は非退化か全等方であり、後者では二平面が等しく $g=3$ になる。
前者では交叉が零で、元の和は五次元である。
\end{proof}
